% 제4장: 실험 결과 및 분석 / 4.1 문체 수정본
% article 계열의 제4절 아래에 삽입하는 본문 조각이다.
% 근거: results/koblex-reproduction-v1/protocol-lock.json,
%       results/koblex-reproduction-v1/chapter4-prep-v1/verification-audit.json,
%       results/koblex-reproduction-v1/chapter4-prep-v1/A1/numeric-ledger.json

\subsection{실험 조건 및 비교 기준}
\label{sec:ch4_conditions}

본 장은 KoBLEX의 공식 한국어 226문항에서 ParSeR의
CE(cross-encoder)를 Jev로 교체했을 때의 근거 선택 성능과 답변 품질을
비교한다. 주 기준선은 공개 코드의 후보 목록 참조 인덱스 문제를 수정한
CE이며, 같은 파이프라인에서 점수화 모델만 Jev로 교체한 조건과
문항별로 대응 비교하였다. 인덱스 문제를 보존한 공개 코드 CE는
재현 점검을 위한 참조 조건으로 두었다(표~\ref{tab:ch4_conditions}).

\begin{table}[htbp]
    \centering
    \caption{비교 조건의 구성}
    \label{tab:ch4_conditions}
    \small
    \setlength{\tabcolsep}{4pt}
    \begin{tabular}{@{}llll@{}}
        \hline
        조건 & 점수화 & 선택 인덱스 & 역할 \\
        \hline
        공개 코드 CE & BGE & 문제 보존 & 재현 참조 \\
        수정 CE      & BGE & 수정      & 주 기준선 \\
        Jev 교체     & Jev & 수정      & 주 비교 \\
        \hline
    \end{tabular}
\end{table}

세 조건은 예상 법령과 BM25 후보를 공유하였다. 예상 법령별 후보
100개를 재정렬하고 상위 10개 중 하나를 공통 LLM 선택기로 선택한 뒤,
선택 근거를 이용해 답변을 생성하였다. 최종 근거 수는 예상 법령 수에
따라 달라지며, 정답 hop 정보는 선택에 사용하지 않았다.
선택ㆍ생성 프롬프트와 평가 기준은 동일하게 유지하였다.
CE는 \texttt{bge-reranker-v2-m3-ko}, Jev는 \texttt{jev-1.14.0}을
사용하였으며, 선택ㆍ답변 생성과 LF-Eval의 평가 모델은
\texttt{gpt-4o-2024-08-06}으로 통일하였다.

정답 근거의 일부가 corpus와 정확히 매핑되지 않는 29문항을 포함하여,
전체 226문항을 주 분석의 분모로 삼았다. 모든 비교는 기존 시험 결과를
확인한 이후 수행한 탐색적 재현 분석이다.
